\section{Conclusion}
\label{sec:conclusion}
This paper presents UL-TransXNet, a lightweight ROI-based TransXNet-family architecture for benign--malignant ultrasound lesion classification. The evidence supports a restrained but coherent conclusion. UL-TransXNet is competitive on TN5000, clearly stronger on BUSI and AUL, and strongest on average across the three public datasets under the unified reporting protocol. TN5000 ablations show that the architecture gains are not explained by a single module: the GG structure is the strongest thyroid-only variant, differential attention improves thresholded behavior, and MCA mainly improves cross-dataset operating robustness. The automatic ROI experiment further shows that detector-predicted crops can connect the classifier to a practical localization front end and remain clearly stronger than full-image classification on BUSI and AUL. Android experiments show that the model can support practical speed--accuracy trade-offs on mobile hardware. At the same time, broad external generalization remains unresolved because all reported datasets are public retrospective benchmarks. Future studies should therefore focus on larger external validation, stronger automatic ROI generation, more rigorous calibration under distribution shift, and deployment evaluation under realistic clinical workflows.
